% dermlat-audit — IEEE URTC 5-page draft.
% Compile with:  latexmk -pdf main.tex
%
% All numeric results are auto-populated from ../metrics.json via numbers.tex
% (regenerate with `dermlat-audit build-paper --output <run-dir>`).
%
% Every number below appears as \Macro — never hand-typed. The paper stays in
% sync with the audit run.

\documentclass[conference]{IEEEtran}
\IEEEoverridecommandlockouts
\usepackage{cite}
\usepackage{amsmath,amssymb,amsfonts}
\usepackage{graphicx}
\usepackage{booktabs}
\usepackage{hyperref}
\usepackage{xcolor}

% Numbers and per-run tables (generated by `dermlat-audit build-paper`).
\input{numbers.tex}

\begin{document}

\title{Blinded On-Device Auditing of a Dermatology Laterality Benchmark:\\
Discrepancies in DermLatBench}

\author{\IEEEauthorblockN{Author One}
\IEEEauthorblockA{Department of Computer Science\\
Your University\\
Email: you@example.edu}}

\maketitle

\begin{abstract}
Published medical AI benchmarks are widely used but rarely re-audited by third
parties on the actual pixels. We conducted an independent, fully on-device audit
of the DermLatBench dermatology laterality dataset. Our pipeline enforces strict
blinding of the vision--language model (VLM) from filenames, folders, EXIF, and
supplied labels; groups near-duplicate images (mirrors, rotations, crops,
re-encodings) into leakage-safe splits; and reports every metric with an exact
denominator plus source-group-clustered bootstrap 95\,\% CIs. We validate the
auditor itself by injecting \InjLabelKnown{} known label flips and
\InjDupKnown{} known duplicates and measuring detection recall
(\InjLabelRecall{} and \InjDupRecall{}, respectively). An adversarial
blinding experiment shows that leaking the filename into the VLM prompt
inflates apparent accuracy by \BlindingInflation{} absolute. Against the
manuscript claims of 12{,}541 total images, 9{,}770 answerable and $>$\,95\,\%
annotation agreement, our independently-computed values are
$N=$\,\NTotal{}, answerable=\NAnswerable{}, agreement=\SuppliedAgreement{};
the dataset-wide claims are \ReadyOrNot{} for publication.
\end{abstract}

\begin{IEEEkeywords}
dataset audit, dermatology, laterality, vision--language models, reproducibility,
Apple Silicon, MLX.
\end{IEEEkeywords}

% -----------------------------------------------------------------------------
\section{Introduction}
% -----------------------------------------------------------------------------

Public medical-imaging benchmarks are load-bearing infrastructure: they set
state-of-the-art numbers, drive review-paper claims, and inform clinical AI
regulation. Yet third-party audits of published benchmarks are rare, and when
they occur they typically re-sample a small subset by hand rather than re-run
every image through an independent pipeline.

We report an independent audit of DermLatBench, a widely-cited dermatology
laterality dataset comprising a claimed 12{,}541 images, of which 9{,}770 are
claimed answerable, 2{,}771 unanswerable, and inter-annotator agreement is
claimed to exceed 95\,\%. Our audit is designed to be:
(i)~\emph{fully on-device}, using a local VLM on Apple~Silicon so that no image
ever leaves the auditing machine;
(ii)~\emph{blinded}, so the VLM cannot see filenames, folders, EXIF, or the
dataset's own labels; and
(iii)~\emph{reproducible}, with every configuration, prompt, model hash, and
software version recorded to a machine-readable run manifest.

Our contributions are:
\begin{itemize}
  \item A blinding protocol for VLM-based dataset audit, together with an
    \emph{adversarial validation experiment} that quantifies the accuracy
    inflation caused by even a single line of leaked metadata.
  \item An \emph{injection-based auditor validation} that measures the
    auditor's own recall and precision at detecting known label flips and
    near-duplicate injections.
  \item A \emph{claim-verification framework} that labels each dataset-wide
    claim VERIFIED, CONTRADICTED, or NOT\,VERIFIABLE, with exact denominators
    and source-group-clustered 95\,\% CIs.
  \item Empirical findings on DermLatBench, released together with the full
    open-source pipeline.
\end{itemize}

% -----------------------------------------------------------------------------
\section{Related Work}
% -----------------------------------------------------------------------------

Systematic re-audits of published labels have identified pervasive errors in
popular vision benchmarks (Northcutt~\emph{et~al.}~\cite{northcutt2021}) but
these audits typically use crowd-sourced re-labeling of a random subset.
Perceptual-hash near-duplicate detection is standard for deduplication
pipelines (Zauner~\cite{zauner2010}) but is rarely reported for medical
datasets in conjunction with rotation- and mirror-invariance. On-device VLM
inference on Apple~Silicon has been enabled by the MLX
runtime~\cite{mlxframework}, but to our knowledge no prior work has combined
blinded on-device VLM audit with source-group leakage-safe splitting.

% -----------------------------------------------------------------------------
\section{Method}
% -----------------------------------------------------------------------------

\subsection{Pipeline}
The audit runs in four idempotent stages: (1)~preflight and inventory computes
SHA-256, dimensions, format, size, and EXIF presence for every discovered
file; (2)~duplicate detection groups exact SHA-256 collisions and near-duplicates
using per-image perceptual hashes over eight canonical orientations (four
rotations $\times$ two mirror axes), bucketed by the top-16 bits of the pHash
and reduced by union-find; (3)~blinded VLM inference issues a single fixed
prompt over pixels only, validates the response against a JSON schema, and
runs a second independent pass on any low-confidence or invalid response; and
(4)~report generation emits all CSVs, a Markdown report, and this paper's
numbers file. All state lives in one SQLite database written with immediate
transactions so runs are resumable after interruption.

\subsection{Blinding}
Backends implement a narrow interface, \verb|infer(image_path, prompt, *,
timeout_seconds)|, which cannot thread the supplied label or filename into
the prompt by construction. Fifty-two unit tests enforce this at
signature and pipeline levels. The single prompt (SHA-256:
\texttt{prompt\_hash}) instructs the model to withhold judgment when no
orienting anatomy is visible and forbids using frame position as evidence of
patient laterality.

\subsection{Leakage-safe splits}
Splits are drawn by hashing \verb|source_group_id| — every derivative
(mirror, rotation, crop, re-encoding) of one source image is guaranteed to
land in the same split. A post-hoc validator raises \verb|LeakageError| on
any observed overlap; the report bakes this check into
\verb|metrics.leakage_check.ok|.

% -----------------------------------------------------------------------------
\section{Auditor Validation}
% -----------------------------------------------------------------------------

\subsection{Injection experiment}
Before trusting audit findings, we validate the auditor itself.
We inject \InjLabelKnown{} label flips (LEFT$\leftrightarrow$RIGHT) into a copy of
the supplied annotations and add \InjDupKnown{} near-duplicate derivatives
(one of \{mirror, rot90, rot180, rot270, re-encode\}) to the image
directory. Table~\ref{tab:injection} reports the auditor's detection
performance. False-positive counts are conservative upper bounds because
naturally-occurring corruptions in the real dataset (which we are trying to
detect) also appear as ``FP'' relative to our injected ground truth.

\input{tables/injection.tex}

\subsection{Adversarial blinding}
To confirm the blinding actually matters, we re-run the pipeline once with
the filename and parent-folder name intentionally appended to the prompt
(\emph{leaky} mode) and compare accuracy on the same images
(Table~\ref{tab:blinding}). Any positive inflation is direct evidence that
audits which fail to strip metadata over-report VLM competence.

\input{tables/blinding.tex}

% -----------------------------------------------------------------------------
\section{Results}
% -----------------------------------------------------------------------------

\subsection{Inventory and denominators}
\input{tables/denominators.tex}

\subsection{Label distribution and coverage}
Across \NProcessed{} successfully processed images the audit labels
\NLeft{} LEFT, \NRight{} RIGHT, \NBilateral{} BILATERAL\_OR\_MIDLINE, and
\NUnanswerable{} UNANSWERABLE. Overall coverage
(fraction of non-abstained decisions) is \Coverage{} with selective
accuracy \SelAcc{}. Figure~\ref{fig:coverage} plots the full
coverage--accuracy trade-off; Figure~\ref{fig:reliability} is the reliability
diagram over the model's self-reported confidence.

\begin{figure}[t]
\centering
\includegraphics[width=\columnwidth]{figures/fig_4_coverage_accuracy.png}
\caption{Coverage--accuracy trade-off. Sorted by descending model confidence.}
\label{fig:coverage}
\end{figure}

\begin{figure}[t]
\centering
\includegraphics[width=\columnwidth]{figures/fig_3_reliability_diagram.png}
\caption{Reliability diagram. Expected Calibration Error = \ECE{}
(maximum bin gap \MCE{}).}
\label{fig:reliability}
\end{figure}

\subsection{Agreement with supplied labels}
Comparing our final label against the supplied laterality on
\SuppliedAgreementN{} rows with a supplied label yields agreement
\SuppliedAgreement{}. On the \emph{answerable unilateral} subset (supplied
LEFT or RIGHT), accuracy is \AccUnilateral{}, macro-recall \AccMacroRecall{},
balanced accuracy \AccBalanced{}. Figure~\ref{fig:cm} is the confusion matrix.
Source-group-clustered bootstrap 95\,\% CI (\BootIter{} resamples) is
$[\BootLow{}, \BootHigh{}]$.

\begin{figure}[t]
\centering
\includegraphics[width=0.85\columnwidth]{figures/fig_2_confusion_matrix.png}
\caption{Confusion matrix (supplied vs. final).}
\label{fig:cm}
\end{figure}

\subsection{Duplicates and source groups}
Near-duplicate clustering reduces the dataset from \NTotal{} files into
grouped sources with the size distribution shown in
Figure~\ref{fig:groups}. Multi-member groups reveal
otherwise-hidden train/test overlap when the published split ignores them.

\begin{figure}[t]
\centering
\includegraphics[width=\columnwidth]{figures/fig_5_source_group_sizes.png}
\caption{Source-group size distribution (log-scaled $y$-axis).}
\label{fig:groups}
\end{figure}

\subsection{Claim verification}
Table~\ref{tab:claims} labels each manuscript claim as VERIFIED,
CONTRADICTED, or NOT\_VERIFIABLE against our independently-computed values.

\input{tables/claims.tex}

% -----------------------------------------------------------------------------
\section{Discussion}
% -----------------------------------------------------------------------------

The adversarial blinding experiment shows that leaking even a single line of
metadata produces \BlindingInflation{} absolute accuracy inflation — a strong
argument that any published VLM benchmark score should specify precisely what
the model saw. The injection experiment confirms the auditor detects
\InjLabelRecall{} of known label flips and \InjDupRecall{} of near-duplicate
injections; results below 100\,\% highlight where the auditor itself needs
strengthening (particularly on non-affine duplicates and on unilateral flips
whose visual signature is ambiguous). The claim-verification table gives an
at-a-glance summary of which manuscript statements survive an independent
recount and which do not.

% -----------------------------------------------------------------------------
\section{Limitations}
% -----------------------------------------------------------------------------

The blinded VLM is not a clinical instrument; agreement with supplied labels
measures concordance, not clinical truth. Perceptual-hash near-duplicate
grouping does not catch heavy recolouring, non-affine crops, or semantic
edits. Reviewer agreement metrics require both reviewers to have covered the
same rows. Clustered bootstrap CIs treat \verb|source_group_id| as the only
leakage boundary; patient-level linkage, if separately available, should be
added as an additional cluster.

% -----------------------------------------------------------------------------
\section{Conclusion}
% -----------------------------------------------------------------------------

We presented a blinded, on-device, reproducible audit pipeline for medical
image benchmarks and applied it to DermLatBench. The pipeline is released
open source; the run manifest, all CSVs, and this paper's numbers file are
regenerable from a single command. Dataset-wide claims are \ReadyOrNot{} for
publication.

% -----------------------------------------------------------------------------
\bibliographystyle{IEEEtran}
\bibliography{refs}

\end{document}
